% Part: lambda-calculus
% Chapter: introduction
% Section: overview

\documentclass[../../../include/open-logic-section]{subfiles}

\begin{document}

\olfileid[hi]{lam}{int}{ovr}
\olsection{रूपरेखा}

लैम्ब्डा कलन मूलतः अलॉन्ज़ो चर्च ने 1930 के दशक के आरंभ में रचनात्मक
तर्कशास्त्र के आधार के रूप में बनाया था, संगणनीय फलनों के मॉडल के रूप में
\emph{नहीं}। किंतु शीघ्र ही यह दिखा दिया गया कि वह संगणनीयता की अन्य
परिभाषाओं, जैसे ट्यूरिंग-संगणनीय फलनों और आंशिक पुनरावर्ती फलनों, के तुल्य
है। आरंभ में इस बात का कुछ आश्चर्यजनक होना इस अभिलक्षण को और भी रोचक बनाता
है।

किसी फलन के लिए कोई नाम परिभाषित करने के बजाय, उसे परिभाषित करने वाले
सांकेतिक व्यंजक से सीधे उसका उल्लेख करने का सुविधाजनक उपाय लैम्ब्डा
संकेत-लिपि है। ``$f$, ~$f(x) = x + 3$ द्वारा परिभाषित फलन हो'' कहने के
बदले हम कह सकते हैं, ``$f$ फलन $\lambd[x][(x + 3)]$ हो।'' दूसरे शब्दों
में, $\lambd[x][(x+3)]$ उस फलन का केवल एक \emph{नाम} है जो अपने तर्क में
तीन जोड़ता है। इस व्यंजक में $x$ एक नाममात्र चर, अथवा स्थानधारक है: उसी
फलन को $\lambd[y][(y + 3)]$ से भी उतनी ही अच्छी तरह दर्शाया जा सकता है। यह
संकेत-लिपि अन्य प्राचलों की उपस्थिति में भी काम करती है। उदाहरणार्थ, मान
लें $g(x, y)$ दो चरों का फलन है और $k$ कोई प्राकृतिक संख्या है। तब
$\lambd[x][g(x,k)]$ वह फलन है जो प्रत्येक~$x$ को~$g(x, k)$ पर प्रतिचित्रित
करता है।

किसी सांकेतिक व्यंजक से फलन परिभाषित करने की इस विधि को \emph{लैम्ब्डा
अमूर्तन} कहते हैं। लैम्ब्डा अमूर्तन का दूसरा पहलू \emph{अनुप्रयोग} है:
मान लें हमारे पास कोई फलन $f$ है (मसलन, प्राकृतिक संख्याओं पर परिभाषित),
तो हम उसे 2 जैसे किसी भी मान पर \emph{लागू} कर सकते हैं। प्रचलित
संकेत-लिपि में परिणाम को निश्चय ही~$f(2)$ लिखते हैं।

लैम्ब्डा अमूर्तन को अनुप्रयोग से जोड़ने पर क्या होता है? तब अमूर्त किए गए
चर के स्थान पर लागू किए जा रहे पद को ``रखकर'' परिणामी व्यंजक को सरल किया
जा सकता है। उदाहरणार्थ,
\[
(\lambd[x][(x + 3)])(2)
\]
को सरल करके~$2 + 3$ किया जा सकता है।

अब तक हमने प्रचलित धारणाओं के लिए नई संकेत-लिपि प्रस्तुत करने के अतिरिक्त
कुछ नहीं किया है। किंतु लैम्ब्डा कलन समुच्चय-सैद्धांतिक दृष्टिकोण से कहीं
अधिक मूलभूत विचलन प्रस्तुत करता है। इस ढाँचे में:
\begin{enumerate}
\item प्रत्येक वस्तु किसी फलन को दर्शाती है।
\item लैम्ब्डा अमूर्तन से फलन परिभाषित किए जा सकते हैं।
\item किसी भी वस्तु को किसी भी दूसरी वस्तु पर लागू किया जा सकता है।
\end{enumerate}
उदाहरणार्थ, यदि $F$ लैम्ब्डा कलन का कोई पद है, तो $F(F)$ को सर्वदा
सार्थक माना जाता है। इस उदार ढाँचे को \emph{टाइपरहित} लैम्ब्डा कलन कहते
हैं, जहाँ ``टाइपरहित'' का अर्थ है, ``किसे किस पर लागू किया जा सकता है, इस
पर कोई प्रतिबंध नहीं।''

\begin{digress}
एक \emph{टाइपयुक्त} लैम्ब्डा कलन भी है, जो टाइपरहित रूपांतर का महत्त्वपूर्ण
भेद है। यद्यपि कई दृष्टियों से टाइपयुक्त लैम्ब्डा कलन टाइपरहित कलन के समान
है, उसे चिरसम्मत समुच्चय-सैद्धांतिक ढाँचे से मिलाना कहीं सरल है और उसके कुछ
गुण बहुत भिन्न हैं।

लैम्ब्डा कलन पर शोध सैद्धांतिक कंप्यूटर विज्ञान और प्रोग्रामिंग भाषाओं के
अभिकल्प में केंद्रीय सिद्ध हुआ है। 1950 के दशक में जॉन मैकार्थी द्वारा
बनाई गई लिस्प ऐसी आरंभिक भाषा का उदाहरण है जिस पर इन विचारों का प्रभाव था।
\end{digress}

\end{document}
